Computational models can reproduce selected measurements from the mouse brain,
but successful prediction does not necessarily reveal the biological mechanism
that generated those measurements. We developed MouseBrainBench to make this
distinction explicit. The framework evaluates reproducibility, predictive
performance, anatomical specificity, temporal direction, and local
structure--function agreement as separate forms of evidence. A mechanistic
interpretation is accepted only when every evidential requirement associated
with that interpretation is satisfied.

We tested this principle in controlled simulations and in three families of
public mouse-brain resources. The framework correctly rejected synthetic
systems that were reproducible but lacked the proposed mechanism. Allen
Neuropixels data produced highly repeatable responses but did not identify a
specific directed topology. Sensorium data supported neural-response prediction
without providing causal evidence. In MICRONS, a local association between
synaptic connectivity and functional readout location survived matched controls
and replicated in two non-overlapping cohorts. This last result agrees with
stronger published MICRONS studies and is used to validate the audit rather than
to claim a new wiring rule.

MouseBrainBench therefore provides a reproducible way to state what partial
digital brain models demonstrate, and equally importantly, what they do not.
